\chapter{Summary and Outlook}
\label{chap:summary}

\section{Achieved Results}
\label{sec:achieved}

The goal of this project was to find out whether a Mixture of Experts model
can learn new data step by step without losing the knowledge it acquired
earlier, and how the router of such a model contributes to this. All five
goals that were fixed in Section~1.3 were reached.

First, the theoretical foundation was worked out. Catastrophic forgetting
happens because every parameter of a dense network takes part in every
task, so the gradient updates for a new task overwrite the solutions of the
earlier ones. The Mixture of Experts architecture attacks exactly this root
cause by routing every input through only a subset of its parameters.

Second, a Mixture of Experts model with a trainable router was designed and
implemented, together with two MNIST based benchmarks and three comparison
approaches, all in a clean and fully reproducible experimental framework.

Third and fourth, the model was trained sequentially on SplitMNIST and
PermutedMNIST and compared quantitatively with a naive baseline, a capacity
matched wide baseline and EWC. On SplitMNIST, the MoE model forgets clearly
less than both baselines. On PermutedMNIST, the first architecture failed,
and the detailed analysis of this failure became one of the most valuable
results of the project: a too strong load balancing loss forces all tasks
to share all experts, and the shared trunk in front of the router stores
exactly the task specific information of this benchmark. An architecture
variant with routing directly at the input, which was derived from this
analysis, reaches the best retention of all compared approaches on
SplitMNIST and recovers a large part of the lost retention on
PermutedMNIST.

Fifth, the router analysis showed at the mechanism level how the protection
works. Different tasks are routed to different expert subsets, so the
gradient updates of a new task mainly rewrite parameters that carry little
old knowledge. The analysis also showed under which conditions this
mechanism fails, namely when the router cannot see the information that
distinguishes the tasks.

Beyond the numbers, the project gave me exactly the practical experience
that was hoped for in the proposal: from the theory of catastrophic
forgetting, through the design and implementation of a non trivial
architecture, to the disciplined experimental evaluation of a research
hypothesis, including an unexpected negative result and its resolution.

\section{Outlook}
\label{sec:outlook}

The results of this work can be extended and transferred in several ways.

\subsection{Extensibility of the Results}

On the architectural side, the model could grow new experts when the router
saturates, which corresponds to the nice to have requirement N2 of
Section~3.3 and would combine the MoE routing with the idea of progressive
networks while keeping the growth sublinear. A gate consistency
regulariser, which punishes changes of the routing distribution on new
data, could stabilise the task to expert assignment further. The load
balancing coefficient, which turned out to be a continual learning
hyperparameter in this project, could be scheduled over time, for example
strong at the beginning of a task and weak afterwards.

A second direction concerns the combination of mechanisms. The evaluation
showed that EWC and MoE protect knowledge in complementary ways: EWC
protects important weights wherever they are, while routing separates
tasks structurally. A routed architecture with a light quadratic penalty
on the shared components, especially on the trunk and the gate itself, is
therefore a natural next step. The input level variant also leaves room
for improvement on PermutedMNIST: its shared hidden layer behind the
experts still mixes all permutations, so a second routed stage could close
the remaining gap to the naive baseline.

\subsection{Transferability of the Results}

On the evaluation side, the approach should be transferred to convolutional
backbones and harder benchmarks such as Split CIFAR-10 or Split
FashionMNIST, and it should be tested in the class incremental scenario,
where the router could additionally serve as a learned task inference
module. The comparison should also be extended to replay based methods,
which were deliberately excluded here. Finally, in the context of very
large pre-trained models, where sparsely gated MoE layers are already
standard, the present results add a small but concrete piece of evidence
that expert routing is not only a capacity scaling device but also a
promising mechanism for lifelong learning, provided the routing is placed
where the tasks actually differ.
